\section{Data limitations}
\begin{itemize}
\item Grades and behaviours are self-reported and were not checked against registrar records.
\item The convenience sample is dominated by KUET, which contributes 79.2\% of the records.
\item The survey is cross-sectional. It measures each student once and cannot observe changes across semesters.
\item CGPA and SGPA are stored as bands rather than consistent exact numbers.
\item Small wording and option differences remain between the two questionnaires after semantic harmonisation.
\item Source A asks for CGPA before the most recent semester, while Source B asks for current CGPA. This one-semester reference difference affects interpretation of the SGPA feature.
\item The dataset does not contain verified course-level grades, longitudinal activity, or equally sized samples from all universities.
\item Questionnaire missing-value modes were calculated before cross-validation and before the fixed train/test split. This historical preprocessing order can transfer limited distribution information into evaluation data.
\end{itemize}
These limitations restrict population-level generalisation and any intended high-stakes use. The model scores are exploratory course-project evidence. They remain useful for reproducible classroom analysis when the sampling, measurement, and preprocessing boundaries stay visible.
